\documentclass{article}
% NeurIPS style, preprint mode (shows authors, no "submitted to" banner) for arXiv.
\usepackage[preprint]{neurips_2023}

\usepackage[utf8]{inputenc}
\usepackage[T1]{fontenc}
\usepackage{hyperref}
\usepackage{url}
\usepackage{booktabs}
\usepackage{amsfonts,amsmath,amssymb}
\usepackage{mathtools}
\usepackage{nicefrac}
\usepackage{microtype}
\usepackage{graphicx}
\usepackage{xcolor}
\usepackage{enumitem}
\usepackage{algorithm}
\usepackage{algpseudocode}
\graphicspath{{figures/}}

% Resilient includes: use a generated artifact if present, else a placeholder.
\newcommand{\inputif}[2]{\IfFileExists{#1}{\input{#1}}{\emph{#2}}}
\newcommand{\figfile}[4]{\IfFileExists{figures/#1}{%
  \begin{figure}[t]\centering\includegraphics[width=#2\linewidth]{#1}%
  \caption{#3}\label{#4}\end{figure}}{}}

\title{An Agentic AI Science Community for\\ Automated Neural Operator Discovery}

\author{%
  Luis Loo \\
  Department of Electrical \& Computer Engineering\\
  Texas A\&M University, College Station, TX\\
  \texttt{luis.loo@tamu.edu} \\
  \And
  Ulisses Braga-Neto \\
  Department of Electrical \& Computer Engineering\\
  Texas A\&M University, College Station, TX\\
  \texttt{ulisses@tamu.edu} \\
}

\begin{document}
\maketitle

\begin{abstract}
Neural operators such as Fourier Neural Operators (FNO) and DeepONets carry
fixed inductive biases, and choosing the right family and architecture for a
given partial differential equation (PDE) remains a manual, expertise-driven
task. We present an \emph{agentic AI Science Community} for automated neural
operator discovery: a swarm of virtual laboratories that collaborate under a
citation-based economy of influence to discover operator architectures with no
human prior on architecture. Each lab couples three role agents---a
\emph{planner} that proposes an architecture, a \emph{worker} that trains and
measures it, and a \emph{reviewer} that participates in cross-lab peer
review---over a shared block vocabulary (Fourier, attention, wavelet, residual
convolution, branch--trunk). The planner and reviewer are driven by a large
language model (LLM); the worker is always numerical (gradient-based training),
by design, so the fitness signal stays comparable across the swarm. We evaluate
on five PDE regimes (piecewise regression, linear advection, Burgers in 1D;
Navier--Stokes and Darcy in 2D), three seeds each, and we \emph{ablate} the LLM
agents against a rule-based particle-swarm coordinator to isolate what
language-model agency contributes. All 9{,}623 LLM calls are logged and audited,
making the claim ``the agents are LLMs'' verifiable rather than assumed: we
observe zero rule-based fallbacks. The ablation yields a mechanistic finding:
at equal training budget the two coordinators reach statistically
indistinguishable screening accuracy, but they discover systematically
\emph{different designs}. The analytic swarm collapses to canonical
single-family answers (pure-Fourier stacks on Navier--Stokes, wavelet stacks on
Darcy), whereas the LLM planner---which chose to hybridize in 99.8\% of its
logged decisions---consistently returns multi-family hybrids that match this
accuracy at lower parameter cost (e.g.\ $2.5\times$ fewer parameters on
Navier--Stokes), and it rediscovers the same four-block advection recipe in
independent seeds. Converged cross-method baselines confirm the honest headline,
matching the operator-learning literature: there is \emph{no universal
winner}---FNO dominates the periodic Navier--Stokes regime while DeepONet-style
models dominate the elliptic Darcy regime---and the converged winners track it:
on Darcy the discovered hybrid attains the best mean error in our baseline
suite, while on Navier--Stokes canonical pure-Fourier stacks remain
unbeaten---the rule-based arm recovers exactly that answer, and the LLM's hybrid
trails it at a matched parameter budget. Each winner degrades sharply on the
other regime. We further
quantify the wall-clock cost of LLM agency. Our contribution is methodological:
reproducible, auditable agentic discovery of competitive hybrid operators. Code,
configurations, and the complete LLM transcripts are released at
\url{https://github.com/luislootx/AI-SC}.
\end{abstract}

\section{Introduction}
Operator learning has become a central tool for building fast surrogates of
partial differential equations (PDEs)~\cite{li2021fno,lu2021deeponet,kovachki2023neuraloperator}.
Yet each operator family encodes a fixed inductive bias: global spectral
convolutions for the Fourier Neural Operator (FNO)~\cite{li2021fno}, a
branch--trunk decomposition for DeepONet~\cite{lu2021deeponet}, multiresolution
bases for wavelet operators~\cite{tripura2023wno}, and attention for
transformer-style operators~\cite{cao2021galerkin,hao2023gnot}. Selecting the
family, depth, width, and block composition for a new PDE is still done by hand,
and---as careful head-to-head studies show---no single family is best across
problem regimes~\cite{lulu2022comparison}. This raises a natural question:

\begin{quote}
\emph{Can a community of automated agents discover the right neural operator
architecture for a given PDE, with no human prior on architecture?}
\end{quote}

We answer this by instantiating the \emph{AI Scientific Community}
framework~\cite{braganeto2026aicommunity}---a swarm of virtual laboratories that explore
a design space under a citation-based economy, where well-regarded labs gain
influence and breeding priority while uncited labs are culled and replaced---for
neural operator discovery. Crucially, we drive the \emph{planning} and
\emph{peer-review} roles of each lab with a large language model, while keeping
the \emph{training} role numerical so that the fitness signal remains comparable
across labs and iterations. This makes the contribution of language-model agency
\emph{measurable}: we ablate the LLM planner/reviewer against an analytic
particle-swarm coordinator under an otherwise identical protocol.

\paragraph{Contributions.}
\begin{enumerate}[leftmargin=1.5em,itemsep=2pt]
\item \textbf{Agentic operator discovery.} An instantiation of the AI Scientific
Community for neural operator architecture search, with explicit LLM planner and
reviewer agents and a deterministic numerical worker
(Section~\ref{sec:framework}--\ref{sec:method}).
\item \textbf{A controlled ablation of LLM agency.} LLM-agent vs.\ rule-based
particle-swarm coordination across five PDE regimes and three seeds, isolating
what the language model adds under an identical training budget. The answer is
architectural rather than accuracy: the LLM returns parameter-lean multi-family
hybrids where the analytic swarm collapses to canonical families
(Section~\ref{sec:results}).
\item \textbf{An evidence-first methodology.} Every LLM call (9{,}623 across the
campaign) is logged to a transcript; we report per-run call counts and fallback
rates, so ``the agents are LLMs'' is auditable, not assumed
(Section~\ref{sec:evidence}, Table~\ref{tab:evidence}).
\item \textbf{A computational-cost analysis.} We quantify per-iteration and
per-lab wall-clock cost, isolating the overhead of LLM agency and situating it
against single-operator training cost at screening fidelity (Section~\ref{sec:cost}).
\item \textbf{An honest cross-regime result.} No operator is universal; the swarm
recovers regime-appropriate designs---on Darcy, a hybrid attaining the best mean
converged error in our baseline suite; on Navier--Stokes, the canonical
pure-Fourier answer via the rule-based arm---and maintains population diversity
rather than collapsing to a single recipe (Section~\ref{sec:results}).
\end{enumerate}

\section{Related Work}
\paragraph{Neural operators.} FNO~\cite{li2021fno} learns a kernel in Fourier
space; DeepONet~\cite{lu2021deeponet} separates input-function encoding (branch)
from query-location encoding (trunk); the neural operator framework unifies these
as learnable maps between function spaces~\cite{kovachki2023neuraloperator}. The
design space has since broadened considerably: wavelet
operators~\cite{tripura2023wno} for multiresolution bases, transformer/attention
operators~\cite{cao2021galerkin,hao2023gnot}, U-shaped multiscale
operators~\cite{rahman2023uno}, and geometry-adapted variants for irregular
domains~\cite{li2023geofno}. Physics-informed
formulations~\cite{raissi2019pinn} and standardized benchmarks~\cite{takamoto2022pdebench}
support systematic comparison. A key empirical lesson, central to our framing, is
that the best operator is problem-dependent: rigorous comparisons find FNO and
DeepONet each win in different regimes~\cite{lulu2022comparison}. The breadth of
this space, and the absence of a universal winner, is precisely what motivates an
automated, per-problem search over operator blocks.

\paragraph{Architecture search.} Neural architecture search (NAS) automates model
design via reinforcement learning~\cite{zoph2017nasrl}, evolutionary
search~\cite{real2019regevo}, gradient-based relaxation of the search
space~\cite{liu2019darts}, and many other strategies~\cite{elsken2019nassurvey}.
A more radical line evolves entire learning algorithms from primitive operations
rather than selecting among predefined blocks~\cite{real2020automlzero}.
Our swarm of labs is an evolutionary/particle-swarm~\cite{kennedy1995pso} search
over an operator-block genome; the novelty here is not the search operator but
that the proposal and selection policies can be delegated to a language model and
ablated against the analytic baseline.

\paragraph{LLMs as optimizers and architecture proposers.} Large language models have
been used as black-box optimizers~\cite{yang2024opro}, as proposal engines for
neural architecture search~\cite{zheng2023genius}, and for program- and
hypothesis-search in mathematical and scientific discovery~\cite{romeraparedes2024funsearch}.
Closest to our planner is EvoPrompting~\cite{chen2023evoprompting}, which uses a
language model as the mutation and crossover operator inside an evolutionary loop
to propose code-level architectures. Our planner plays the same role over an
operator-block genome, but we make two design choices that sharpen measurability:
the fitness signal comes from a deterministic training procedure rather than a
learned predictor, and every LLM-driven run is paired with a rule-based
particle-swarm~\cite{kennedy1995pso} arm so the language model's contribution is
isolable by ablation.

\paragraph{Multi-agent LLM systems.} A growing body of work composes multiple
LLM agents with distinct roles that communicate to solve a task: conversational
multi-agent frameworks~\cite{wu2024autogen}, communicative agent
societies~\cite{li2023camel}, and role-specialized teams that emulate a software
company~\cite{qian2024chatdev} or a standardized-operating-procedure
organization~\cite{hong2024metagpt}. At the workflow level such teams drive
end-to-end research agents~\cite{lu2026aiscientist} and numerical algorithm
design~\cite{toscano2025athena}. Our framework instantiates this paradigm with a
planner--worker--reviewer triad per lab; the reviewer's accept/reject vote is a
form of LLM-as-judge evaluation~\cite{zheng2023llmjudge}, here applied to peer
review of candidate operators rather than to chatbot responses. We build directly
on the AI Scientific Community model of agentic virtual-lab
swarms~\cite{braganeto2026aicommunity}, and contribute a controlled,
ablation-backed study of when the language-model agents actually outperform their
rule-based counterparts.

\section{Problem Formulation}\label{sec:problem}
\paragraph{Operator learning.} Given a PDE with solution operator
$\mathcal{G}:\mathcal{A}\to\mathcal{U}$ mapping an input function $a\in\mathcal{A}$
(e.g.\ an initial condition or coefficient field) to the solution
$u=\mathcal{G}(a)\in\mathcal{U}$, we seek a parametric surrogate
$\mathcal{G}_\theta$ minimizing the expected relative $L_2$ error
\begin{equation}
\mathcal{L}(\theta) \;=\; \mathbb{E}_{a\sim\mu}
\frac{\lVert \mathcal{G}_\theta(a) - \mathcal{G}(a)\rVert_2}
     {\lVert \mathcal{G}(a)\rVert_2},
\end{equation}
estimated on sampled input--output pairs. The relative $L_2$ error is our primary
accuracy metric throughout.

\paragraph{Architecture search space.} We define an architecture by a
\emph{genome} $g=(b_{1:L},\,c,\,m,\,\sigma,\,\dots)$, where $b_{1:L}$ is a
sequence of $L\in[2,8]$ blocks drawn from a vocabulary
$\mathcal{B}=\{\textsc{fourier},\textsc{attention},\textsc{wavelet},
\textsc{residual\_conv},\textsc{branch\_trunk}\}$, $c$ is the hidden width, $m$
the number of spectral modes, $\sigma$ the activation, plus gating, skip
connections, dropout, and learning rate. A genome compiles to a concrete neural
operator $\mathcal{G}_{\theta(g)}$. The discovery problem is the bilevel program
\begin{equation}
g^\star \;=\; \arg\min_{g}\; \mathcal{L}\big(\theta^\star(g)\big),
\qquad
\theta^\star(g) \;=\; \arg\min_{\theta}\; \widehat{\mathcal{L}}_{\text{train}}\big(\theta; g\big),
\end{equation}
where the inner problem is ordinary gradient-based training of a fixed
architecture and the outer problem searches genomes. We solve the outer problem
with a swarm of agentic labs (Section~\ref{sec:framework}); the inner problem is
always numerical.

\section{The AI Scientific Community Framework}\label{sec:framework}
A swarm of $N$ virtual labs explores the genome space. Each lab $i$ holds a
current genome $g_i$, a personal best, a velocity (for the analytic update), and
a \emph{trust} score accumulated through peer review. Each iteration runs four
phases (Algorithm~\ref{alg:swarm}, Figure~\ref{fig:method}):

\begin{enumerate}[leftmargin=1.5em,itemsep=1pt]
\item \textbf{Planning.} Each lab proposes its next genome $g_i$.
\item \textbf{Training.} Each lab compiles $g_i$ and trains it at screening
fidelity, returning measured relative $L_2$ (clean and noisy).
\item \textbf{Peer review.} Labs review a sample of peers and cast soft votes;
votes accumulate as citation-like influence and update trust scores.
\item \textbf{Lifecycle.} A composite fitness combines accuracy, generalization,
parameter efficiency, and novelty (with an accuracy floor); the global best is
updated, influential labs persist and breed, and uncited labs are culled and
replaced, realizing selection over research directions.
\end{enumerate}

The composite fitness for lab $i$ is
\begin{equation}
F_i = w_a\,A_i + w_g\,G_i + w_e\,E_i + w_n\,\nu_i,
\qquad (w_a,w_g,w_e,w_n)=(0.70,0.20,0.05,0.05),
\end{equation}
zeroed when accuracy falls below a floor, which prevents trivially small or
degenerate models from winning on efficiency or novelty alone. Here $A_i$ is an
accuracy score derived from relative $L_2$, $G_i$ rewards robustness to input
noise, $E_i$ rewards parameter parsimony, and $\nu_i$ is an architectural-novelty
score relative to the rest of the population. The weights are fixed a priori,
shared by every run and both conditions, and were not tuned; accuracy dominates
by design, and ablating the weighting itself is left to future work.

\begin{algorithm}[t]
\caption{AI Scientific Community swarm (one configuration)}
\label{alg:swarm}
\begin{algorithmic}[1]
\State initialize $N$ labs across paradigms (FNO, DeepONet, transformer, wavelet, random/hybrid)
\For{iteration $t=1,\dots,T$}
  \For{each active lab $i$} \Comment{Phase 1: plan}
     \State $g_i \gets \textsc{Plan}(g_i,\, g^{\text{best}},\, \{g_j\}_{j\neq i},\, \text{history}_i)$
     \Comment{LLM agent \emph{or} PSO update}
  \EndFor
  \For{each active lab $i$} \Comment{Phase 2: train (always numerical)}
     \State $\theta_i \gets \textsc{Train}(g_i)$; \; measure relative $L_2$
  \EndFor
  \State $\{v_i\} \gets \textsc{PeerReview}(\{(g_i,\text{fitness}_i)\})$
  \Comment{Phase 3: LLM agent \emph{or} accuracy-proportional}
  \State update composite fitness, trust, global best; cull uncited, breed influential
  \Comment{Phase 4}
  \State checkpoint state \Comment{resume-safe}
\EndFor
\end{algorithmic}
\end{algorithm}

\figfile{fig_method2.png}{1.0}{The agentic AI Science Community. Each of five
PDE regimes (left) is attacked by a community of $N=16$ virtual labs (center,
schematic): inside every lab, a \emph{planner} (LLM) proposes an operator-block
genome $g_i$, a deterministic \emph{worker} trains it (PyTorch) and measures
relative $L_2$, and a \emph{reviewer} (LLM) casts soft peer-review votes over
peer candidates. Votes act as citations: influential labs gain trust and breed,
uncited labs are culled and replaced, over $t=1\dots20$ iterations. The two
decision roles are ablated against analytic PSO rules under an identical
training budget. The community returns a different architecture per regime
(right): the campaign's converged winners of Table~\ref{tab:baselines}.}{fig:method}

\section{Method: Agentic Roles}\label{sec:method}
\subsection{Three role agents per lab}
Each lab couples three roles. The \textbf{planner} proposes the next genome; the
\textbf{worker} compiles and trains the corresponding operator with gradient
descent, returning measured relative $L_2$; the \textbf{reviewer} casts
peer-review votes over a sample of peer candidates. The \textbf{worker is always
numerical}: we deliberately do not delegate training to a language model, so the
fitness signal is consistent and comparable across labs and iterations. Thus
``LLM agents'' in this paper refers precisely to the \emph{planner} and
\emph{reviewer} roles---the decision-makers---never to operator training.

\subsection{LLM-agent and rule-based modes}
The framework is model-agnostic in the planner/reviewer roles, which lets us run
a clean ablation:
\begin{itemize}[leftmargin=1.5em,itemsep=2pt]
\item \textbf{LLM agents.} The planner queries an LLM for the next genome, given
its own history, the global best, and a summary of peer labs and their measured
fitness; the reviewer queries an LLM to score candidate genomes against measured
fitness. Proposals are parsed from JSON and snapped to the valid genome space. We
serve \texttt{gemma3:12b} locally via Ollama, so the pipeline is reproducible
with no paid API.\footnote{We selected this model on a structured-output
criterion rather than raw scale. A smaller efficiency-tier variant
(\texttt{gemma4:e4b}) failed to emit schema-valid JSON for the reviewer
role---even with prompt hardening and an enlarged context window---and silently
degraded to a non-LLM fallback; a replay probe over real planner and reviewer
prompts confirmed that \texttt{gemma3:12b} produces valid structured output for
both roles where the smaller variant did not. The evidence audit
(Section~\ref{sec:evidence}) exists precisely to make such silent degradation
detectable, and verifies that none occurs in the reported runs.}
\item \textbf{PSO (rules).} The planner uses analytic explore/exploit
particle-swarm updates~\cite{kennedy1995pso}; the reviewer uses
accuracy-proportional voting. No language model is involved.
\end{itemize}
Both modes share swarm size, iteration count, training budget, fitness, and
seeds; only the planner/reviewer policy differs. This isolates the effect of
language-model agency.

\subsection{Evidence that the agents are LLMs}\label{sec:evidence}
A recurring risk in ``agentic'' systems is that a language model silently
degrades to a rule-based fallback (e.g.\ on a transient error or malformed
output) while still being described as an agent. We guard against this by logging
\emph{every} LLM call---system prompt, user prompt, raw response, role, and
parse status---to a per-run transcript (\texttt{llm\_transcript.jsonl}), and by
auditing each run's log for planner/reviewer fallbacks. Table~\ref{tab:evidence}
reports, per LLM run, the planner and reviewer call counts, the parse-success
rate, and the fallback count. A fallback count of zero certifies that planning
and review were LLM-driven throughout; the transcripts are released so this audit
is independently reproducible.

\subsection{Screening fidelity}
Each candidate is trained at deliberately reduced ``screening'' fidelity (small
data, few epochs, modest resolution) so a single operator trains in seconds, not
hours, making a swarm of hundreds of trainings tractable on one GPU
(Section~\ref{sec:cost}). Absolute screening errors are therefore higher than
fully-converged training; we report converged cross-method baselines separately
to calibrate the discovered architectures.

\section{Experimental Setup}\label{sec:setup}
We run five PDE regimes: piecewise regression, linear advection, and Burgers in
1D; Navier--Stokes (vorticity form, Kolmogorov forcing) and Darcy flow in 2D.
Each configuration runs for three seeds ($42,137,2024$) in both the LLM and PSO
conditions ($5\times2\times3=30$ runs), with $N=16$ labs for $T=20$ iterations
(the paper-grade preset). All runs are resume-safe (checkpointed per iteration)
to survive interruptions. We additionally train converged cross-method baselines
(FNO at several sizes, DeepONet, POD-DeepONet, a pure transformer, and the
swarm-discovered hybrid) with input/output normalization, three seeds each. All
timings are on a single NVIDIA RTX~4080.

\section{Results}\label{sec:results}
\paragraph{No universal winner.}
Table~\ref{tab:baselines} reports converged
cross-method baselines on Navier--Stokes and Darcy. The pattern matches the
operator-learning literature~\cite{lulu2022comparison} and is the honest core of
our story: FNO dominates the periodic, spectral-friendly Navier--Stokes regime by
an order of magnitude, while DeepONet-style models dominate the elliptic,
Dirichlet Darcy regime. No single operator wins both. The same table evaluates
the swarm's two discovered winners, retrained from scratch to convergence, on
\emph{both} regimes. Each is competitive with the regime-best family at home.
The Darcy hybrid attains $0.0631\pm0.0083$, the best mean Darcy error in the
suite, statistically tied with DeepONet faithful ($0.0649\pm0.0037$) at three
seeds. The NS hybrid reaches $0.0005$ at $1.1$M parameters---an order of
magnitude below every non-FNO baseline and $30\times$ fewer parameters than the
large FNO---but we note plainly that a canonical pure-Fourier stack at the
\emph{same} budget is better still (FNO h32: $0.0002$ at $1.2$M), and that this
canonical stack is exactly what the PSO arm discovers on NS. Off-regime, each
winner degrades: the NS hybrid collapses on Darcy ($0.71$) and the Darcy hybrid
gives up $3.6\times$ on NS ($0.0018$). The pattern is symmetric and honest:
rule-based search wins precisely where the textbook answer is optimal, LLM
hybridization pays where no single family suffices, and neither
transfers---exactly the no-free-lunch behavior that motivates discovering a
different architecture per regime rather than reusing one.

\begin{table}[t]\centering
\caption{Converged cross-method baselines (relative $L_2$, mean over seeds;
parameters counted at this protocol's resolution---block parameter counts scale
with grid size). FNO wins Navier--Stokes; DeepONet-style models win Darcy. The
two ``Discovered hybrid'' rows are the campaign's winning architectures per
regime, each retrained from scratch and evaluated on \emph{both} regimes: each
is competitive at home and degrades off-regime, which is why per-regime
discovery is needed.}
\label{tab:baselines}
\inputif{tables/tab_baselines.tex}{(baselines table is generated by the build pipeline.)}
\end{table}

\paragraph{Agentic discovery: hybrids per regime vs.\ canonical families.}
Table~\ref{tab:discovery} reports, per PDE, the relative $L_2$ and parameter
count of the architecture discovered under the LLM and PSO conditions;
Figure~\ref{fig:fig_archdiagram.pdf} visualizes the discovered block sequences,
and Appendix~\ref{app:perseed} lists every per-seed winner. Two patterns emerge.
First, \emph{every} LLM winner is a multi-family hybrid, and the mixture is
regime-appropriate: Fourier-led stacks on the 1D transport problems (on
advection, two of three seeds independently discover the identical four-block
recipe \texttt{F-A-W-T}: fourier, attention, wavelet, branch--trunk), and
attention/wavelet-led mixtures on piecewise regression and Darcy. Second, the
PSO condition instead collapses to canonical, nearly single-family designs:
pure-Fourier stacks on Navier--Stokes (two of three seeds---the textbook answer
for a periodic, spectrally friendly regime) and wavelet-dominated stacks on
Darcy and piecewise regression. The trade-off against established baselines is
shown in Figure~\ref{fig:fig_pareto.pdf}: the discovered hybrids sit on a
competitive accuracy--parameter frontier; on Navier--Stokes the LLM hybrids
match the pure-spectral PSO accuracy with $2.5\times$ fewer parameters on
average ($2.2$M vs.\ $5.6$M).

\begin{table}[t]\centering
\caption{Swarm-discovered architectures: relative $L_2$ (clean) and parameter
count, LLM (Gemma~3) vs.\ PSO, mean$\pm$std over three seeds.}
\label{tab:discovery}
\inputif{tables/tab_discovery.tex}{(discovery table populates as the campaign completes.)}
\end{table}

\figfile{fig_archdiagram.pdf}{0.85}{Discovered architecture per PDE (LLM agents,
seed 42): every regime's winner is a multi-family hybrid---Fourier-led on the 1D
transport problems, attention/wavelet-led on piecewise regression and
Darcy---with a different recipe per regime.}{fig:fig_archdiagram.pdf}

\figfile{fig_pareto.pdf}{1.0}{Accuracy--parameter trade-off on the 2D problems.
Swarm labs (circles, colored by paradigm) and converged baselines (stars). The
swarm populates a competitive frontier at far fewer parameters than the large
FNO baselines.}{fig:fig_pareto.pdf}

\paragraph{Ablation: what does LLM agency add?}
Figure~\ref{fig:fig_ablation.pdf} compares discovered relative $L_2$ under LLM and
PSO coordination, and Figure~\ref{fig:fig_convergence.pdf} compares convergence of
the swarm's best error per iteration. Because both conditions train identical
operators under an identical budget, any difference is attributable to the
planner/reviewer policy. At three seeds the screening accuracies are
statistically indistinguishable on all five PDEs (overlapping $\pm 1$ s.d.), and
we do not claim an accuracy win for LLM agency. The measurable difference is
architectural: the LLM discovers parameter-lean multi-family hybrids where PSO
converges to canonical families, and it does so reproducibly (the identical
advection recipe across independent seeds). We report this plainly rather than
claiming LLM dominance: the value of the framework is the agentic, auditable
discovery process, the ablation that measures it, and the design-space behavior
it reveals.

\paragraph{How the LLM plans.}
The transcripts let us inspect the planner's \emph{decisions}, not only its
outcomes. Across 4{,}817 planner calls the LLM chose \textsc{hybridize} in
99.8\% of decisions (the remainder \textsc{explore}; never \textsc{exploit}),
and its rationales condition on the summarized swarm state, e.g.\ ``\emph{several
labs are converging on FNO and transformer architectures\ldots\ I'll introduce a
hybrid approach by incorporating attention blocks}'' (Navier--Stokes, seed 42).
This action bias is the mechanism behind the architectural contrast above: a
planner that almost always grafts blocks from other families structurally
produces multi-family hybrids, whereas the analytic update can---and
does---collapse to single-family optima. The bias is partly prompt-induced (the
planner is instructed to preserve diversity when the swarm converges;
Appendix~\ref{app:prompts}), so we read it as a caveat as much as a capability:
LLM agency here manifests as context-conditioned hybridization rather than
balanced explore/exploit arbitration.

\paragraph{What the reviewer rewards.}
The same transcripts quantify the reviewer. Across 4{,}804 scoreable review
calls, the rank correlation between the reviewer's vote distribution and the
candidates' measured accuracy is $\bar{\rho}=0.64$ (Spearman; positive in 88\%
of calls, a perfect accuracy ordering in 50\%), so peer review is grounded in
the fitness signal rather than stylistic preference. Its rationales cite
accuracy in 97\% of calls, parameter efficiency in 83\%, and robustness in
71\%, but novelty in only 26\%: the deviations from pure accuracy ordering are
consistent with the prompt's instruction to also reward novelty and efficiency.
Review is thus fitness-anchored with a diversity tilt---precisely the property
the citation lifecycle needs, since accurate labs must accumulate influence
without unusual designs being starved (Figure~\ref{fig:fig_blockevo.pdf}).

\figfile{fig_ablation.pdf}{0.7}{Discovered relative $L_2$ per PDE: LLM agents
(Gemma~3) vs.\ rule-based PSO, under an identical training budget.}{fig:fig_ablation.pdf}

\figfile{fig_convergence.pdf}{1.0}{Swarm convergence: best relative $L_2$ per
iteration, LLM vs.\ PSO, mean $\pm$ s.d.\ over seeds.}{fig:fig_convergence.pdf}

\paragraph{The swarm maintains diversity.}
A failure mode of citation-driven swarms is premature collapse onto an early
leader (groupthink). Figure~\ref{fig:fig_blockevo.pdf} shows population-level
block usage stays mixed across iterations: the novelty term and minimum
exploration keep multiple paradigms alive even as individual winners specialize.
The ``different recipe per regime'' result (Figure~\ref{fig:fig_archdiagram.pdf})
is thus a property of the \emph{winner}, not a collapse of the population.

\figfile{fig_blockevo.pdf}{1.0}{Population block-usage fraction across iterations
(LLM). Diversity is maintained---no single block type dominates---indicating the
swarm avoids groupthink collapse.}{fig:fig_blockevo.pdf}

\paragraph{The agents really were LLMs.}
Table~\ref{tab:evidence} summarizes the transcript audit over all LLM runs:
planner and reviewer calls, parse-success, and fallbacks. Across the campaign we
observe zero fallbacks, certifying that planning and review were LLM-driven.

\begin{table}[t]\centering
\caption{Evidence audit: per-run planner/reviewer LLM call counts, parse-success,
and fallbacks. Zero fallbacks certifies LLM-driven planning and review. Counts
above the nominal $16\times20=320$ reflect re-planned iterations after mid-run
interruptions (the runs are checkpointed and resume-safe); every call, including
repeats, is in the released transcripts.}
\label{tab:evidence}
\inputif{tables/tab_evidence.tex}{(evidence table populates as the campaign completes.)}
\end{table}

\section{Computational Cost}\label{sec:cost}
Table~\ref{tab:timing} and Figure~\ref{fig:fig_timing.pdf} report wall-clock times
on a single RTX~4080: the total time for a full $16$-lab $\times$ $20$-iteration
discovery run, and the time per lab-evaluation (one lab's plan--train--review
cycle), for both conditions. Because the LLM and PSO conditions train identical
operators under an identical budget, the difference in per-lab time is precisely
the overhead of LLM-driven planning and review (two language-model calls per lab
per iteration): roughly $14$s per lab-evaluation on the 1D problems and
$25$--$33$s on the 2D problems, i.e.\ a $2$--$4\times$ increase in total
wall-clock over rule-based coordination. The LLM condition also carries higher
run-to-run variance on the 2D problems (e.g.\ $271\pm104$ min on Navier--Stokes):
an honest consequence of letting the planner commit to parameter-heavy hybrids
whose training dominates the budget. The overhead buys the behavior of
Section~\ref{sec:results}; whether that trade is worthwhile is regime- and
goal-dependent, and the ablation makes it measurable.

\begin{table}[t]\centering
\caption{Computational cost on a single RTX~4080: total run time and per
lab-evaluation time, LLM (Gemma~3) vs.\ PSO, mean$\pm$std over seeds. The LLM--PSO
gap isolates the cost of language-model agency.}
\label{tab:timing}
\inputif{tables/tab_timing.tex}{(timing table populates as the campaign completes.)}
\end{table}

\figfile{fig_timing.pdf}{0.7}{Per lab-evaluation wall time, LLM vs.\ PSO. The gap
is the cost of LLM planning and peer review at identical training budget.}{fig:fig_timing.pdf}

\paragraph{Why a swarm of hundreds of trainings is feasible.}
The swarm screens architectures at reduced fidelity, so a single operator trains
in seconds, not hours. Table~\ref{tab:trainop} lists single-operator training
times to convergence at screening fidelity on Navier--Stokes; cost is dominated
by epochs and block type, not parameter count (a 33.6M-parameter FNO trains
faster than a 105K-parameter transformer because attention is quadratic in
sequence length). A full discovery run evaluates $16\times20=320$ architectures
within the budgets above, reserving expensive full-fidelity training for the
final discovered architecture and the baseline comparison---the standard
multi-fidelity argument: screen cheaply, commit expensively.

\begin{table}[t]\centering
\caption{Single-operator training time to convergence at screening fidelity
(Navier--Stokes), mean over seeds.}
\label{tab:trainop}
\inputif{tables/tab_trainop.tex}{(training-time table is generated by the build pipeline.)}
\end{table}

\section{Discussion and Limitations}
\paragraph{LLM-as-policy, not autonomous tool use.} Our agents use the LLM as a
\emph{policy} for the planning and review decisions; the LLM does not execute code
or call external tools. This is a narrower notion of ``agentic'' than fully
autonomous research agents~\cite{lu2026aiscientist}, and we state it plainly.
\paragraph{Joint role ablation.} The ablation switches planner and reviewer
together. The transcript analyses provide observational evidence about each
role's behavior, but disentangling their causal contributions requires a
factorial design (LLM planner with analytic reviewer, and vice versa), which we
leave to future work.
\paragraph{Literature priors.} A mid-size or frontier LLM has read the
operator-learning literature and may ``know'' that FNO suits periodic problems.
This is simultaneously a feature (informed priors) and a confound (it is not
prior-free discovery); the PSO condition provides the prior-free reference point.
\paragraph{Determinism.} LLM planning is stochastic; we report multiple seeds and
release transcripts rather than claiming bit-for-bit reproducibility.
\paragraph{Screening fidelity.} The swarm screens at reduced fidelity; absolute
errors should be read against the converged baselines, not as final accuracies.
\paragraph{Model scale.} We use a single local model (\texttt{gemma3:12b});
whether a frontier model changes discovery quality is left to future work, enabled
by the model-agnostic backend.

\section{Reproducibility}\label{sec:repro}
Code, the resume-safe runner, exact configurations, and the complete per-run LLM
transcripts are available at \url{https://github.com/luislootx/AI-SC}. Every
reported LLM run ships its \texttt{llm\_transcript.jsonl}; the evidence audit
(Table~\ref{tab:evidence}) is reproducible directly from these files. The LLM
backend (\texttt{gemma3:12b} via Ollama) runs locally with no paid API. Representative
planner and reviewer prompts are listed in Appendix~\ref{app:prompts}.

\section{Conclusion}
We presented an agentic AI Science Community that discovers neural operator
architectures through LLM-driven planning and peer review over a swarm of virtual
labs, with a deterministic numerical worker and a controlled ablation against
rule-based coordination. The honest cross-regime finding is that no operator is
universal; the contribution is a reproducible, auditable methodology for agentic
architecture discovery that recovers regime-appropriate, parameter-lean hybrids
at a quantified LLM overhead. Future work includes frontier LLM backends, a
factorial planner/reviewer ablation, larger block vocabularies (graph operators,
neural Galerkin), action-space debiasing of the planner, and full-fidelity
re-training of every discovered architecture.

\begin{ack}
We thank the organizers of the ICERM Hot Topics Workshop on Agentic Scientific
Computing and Scientific Machine Learning, where a preliminary version of this
work was presented as a poster, and the Texas A\&M ECEN graduate program for
travel support.
\end{ack}

\bibliographystyle{plain}
\bibliography{refs}

\appendix
\section{Agent Prompts}\label{app:prompts}
The planner and reviewer system prompts are reproduced verbatim below (PDE name
instantiated per run); each user prompt carries the lab's current genome, its
recent fitness records, the global best, and an anonymized summary of peer labs,
exactly as recorded in every run's \texttt{llm\_transcript.jsonl}.

\paragraph{Planner system prompt.}
{\small\begin{verbatim}
You are the planning agent of a virtual research lab in a
swarm of labs collectively discovering neural operator architectures for 2D
Navier-Stokes. Each lab is one PSO particle; your job is to propose the next
architecture (genome) for your lab to train and evaluate next iteration.

You must balance EXPLORATION (novel hybrids, unexplored block combinations)
with EXPLOITATION (move toward the global best architecture observed in the
swarm). Use the lab's current exploration_rate as a guide: high -> favor
exploration, low -> favor exploitation. Consider what other labs are doing --
if the swarm is converging, propose something different to maintain diversity
and avoid groupthink.

Output STRICT JSON ONLY (no prose, no fences) matching this schema:
{
  "action": "explore" | "exploit" | "hybridize",
  "rationale": "<1-2 sentences explaining your choice>",
  "new_genome": { "block_sequence": [...], "hidden_channels": <int>,
    "fourier_modes": <int>, "activation": "<str>", "use_gating": <bool>,
    "use_skip_connections": <bool>, "dropout_rate": <float>,
    "learning_rate": <float>, "weight_decay": <float> }
}
- ALLOWED_BLOCKS = ["fourier","attention","branch_trunk","wavelet",
                    "residual_conv"]
\end{verbatim}}

\paragraph{Reviewer system prompt.}
{\small\begin{verbatim}
You are an anonymous peer reviewer evaluating neural
operator architectures proposed by other virtual labs in a research swarm.

You must rank the candidates on overall scientific merit for solving 2D
Navier-Stokes operator learning. Reward: high accuracy, good generalization
under noise, parameter efficiency, architectural novelty (unusual block
combinations), and theoretical soundness (e.g. fourier blocks for periodic
domains, attention for nonlocal coupling). Penalize: groupthink (identical to
global best), pathological setups (too few or too many blocks, mismatched
modes), or poor measured fitness.

Output STRICT JSON ONLY:
{
  "scores":   { "<lab_id>": <float in [0,1]>, ... },
  "rationale": "<1-2 sentences total, not per lab>"
}
Scores must sum to roughly 1.0 across the candidates (soft-vote distribution).
\end{verbatim}}

\section{Per-seed Results}\label{app:perseed}
Table~\ref{tab:perseed} lists the discovered block sequence, screening relative
$L_2$, and parameter count for every run of the campaign (5 PDEs $\times$ 2
conditions $\times$ 3 seeds). Machine-readable versions are generated into
\texttt{results/aggregate\_v2.json} in the repository.

\begin{table}[h]\centering
\caption{Per-seed discovered architectures, all 30 runs. Blocks abbreviated
F=fourier, A=attention, W=wavelet, R=residual\_conv, T=branch\_trunk.}
\label{tab:perseed}
{\small
\inputif{tables/tab_perseed.tex}{(per-seed table is generated by the build pipeline.)}
}
\end{table}

\section{Experimental Details}\label{app:details}
Table~\ref{tab:expdetails} lists the screening-phase configuration per PDE
(from each run's released \texttt{config.json}). Both conditions share these
budgets exactly; only the planner/reviewer policy differs. Every candidate is
trained from scratch for the stated epochs each iteration; robustness
(noisy relative $L_2$) is measured under additive input noise
$\varepsilon\sim\mathcal{N}(0,\,0.1^2)$. The swarm uses $N=16$ labs for $T=20$
iterations in all runs. The LLM backend is \texttt{gemma3:12b} served by Ollama
with an 8192-token context window; all timings are on a single RTX~4080.
Converged-baseline budgets (Table~\ref{tab:baselines}) are per-family---FNO
variants 50 epochs, hybrids and POD-DeepONet and the pure transformer 80,
DeepONet-faithful 600 with input/output normalization---as encoded in the
released \texttt{validate\_baselines\_v3.py}.

\begin{table}[h]\centering
\caption{Screening-phase configuration per PDE (identical across conditions and
seeds).}
\label{tab:expdetails}
\begin{tabular}{lccccc}
\toprule
PDE & Dim & Resolution & Train / test & Epochs per candidate \\
\midrule
Piecewise Reg. & 1D & 128 & 256 / 64 & 40 \\
Advection & 1D & 128 & 256 / 64 & 40 \\
Burgers & 1D & 128 & 256 / 64 & 40 \\
Navier--Stokes & 2D & $32\times32$ & 512 / 128 & 50 \\
Darcy & 2D & $32\times32$ & 512 / 128 & 50 \\
\bottomrule
\end{tabular}
\end{table}

\end{document}
